\documentclass[10pt]{scrartcl}
\usepackage[mdthm,colorsec,noauthor,noasy]{evan}
\usepackage{handout}
\DeclareMathOperator{\spn}{span}\DeclareMathOperator{\nul}{null}\DeclareMathOperator{\range}{range}
\newcommand{\PP}{\mathcal P}\newcommand{\MM}{\mathcal M}
\title{Math 110 --- Week 4}
\subtitle{Solutions}
\author{}\date{}\ihead{}
\begin{document}
\maketitle

\begin{prob}\textbf{HW4 P1.}\par\idea{The identity in (a) hands you the inverse when $ad-bc\ne0$, and forces $A=0$ when $ad-bc=0$ and $A$ is assumed invertible.}
\begin{steps}
  \ii (a) Write $B=\begin{psmallmatrix}d&-b\\-c&a\end{psmallmatrix}$. $AB=\begin{psmallmatrix}ad-bc&-ab+ba\\cd-dc&-cb+da\end{psmallmatrix}=(ad-bc)I$. Likewise $BA=\begin{psmallmatrix}da-bc&db-bd\\-ca+ac&-cb+ad\end{psmallmatrix}=(ad-bc)I$.
  \ii (b) If $\delta=ad-bc\ne0$, then $A(\delta^{-1}B)=I=(\delta^{-1}B)A$. So $A$ is invertible with $A^{-1}=\frac1{ad-bc}\begin{psmallmatrix}d&-b\\-c&a\end{psmallmatrix}$.
  \ii (c) Suppose $ad-bc=0$ and $A$ is invertible. Then $AB=0$, and multiplying on the left by $A^{-1}$ gives $B=0$. So $d=b=c=a=0$, i.e.\ $A=0$. But the zero matrix is not invertible. Contradiction.
\end{steps}
\end{prob}

\begin{prob}\textbf{HW4 P2.}\par\idea{Compute $ad-bc$ in the stated field.}
\begin{steps}
  \ii $\begin{psmallmatrix}2&1\\5&3\end{psmallmatrix}$: $6-5=1$. Invertible, inverse $\begin{psmallmatrix}3&-1\\-5&2\end{psmallmatrix}$.
  \ii $\begin{psmallmatrix}1&i\\i&-1\end{psmallmatrix}$: $(1)(-1)-i\cdot i=-1+1=0$. Not invertible. (Second column is $i$ times the first.)
  \ii $\begin{psmallmatrix}1&2\\3&1\end{psmallmatrix}$ over $\FF_5$: $1-6=-5=0$ in $\FF_5$. Not invertible. (Second column is $2$ times the first: $2\cdot3=6=1$.)
\end{steps}
\end{prob}

\begin{prob}\textbf{HW4 P3.}\par\idea{Injective maps keep independence; $n$ independent vectors in dimension $n$ are a basis. Conversely a basis in the range makes $T$ surjective.}
\begin{steps}
  \ii (a)$\Rightarrow$(b). Let $v_1,\dots,v_n$ be a basis. If $\sum\lambda_jTv_j=0$ then $T(\sum\lambda_jv_j)=0$, so $\sum\lambda_jv_j=0$ ($T$ injective), so all $\lambda_j=0$.
  \ii So $Tv_1,\dots,Tv_n$ are $n$ independent vectors in $V$, $\dim V=n$: a basis.
  \ii (b)$\Rightarrow$(a). Take any basis $v_1,\dots,v_n$. Then $\range T\supseteq\spn(Tv_1,\dots,Tv_n)=V$, so $T$ is surjective, hence an isomorphism (Prop.~31).
\end{steps}
\end{prob}

\begin{prob}\textbf{HW4 P4.}\par\idea{Rank is the dimension of the range. Compare the range of $BC$ with the range of $B$, and view it as the image of $\range C$.}
\begin{steps}
  \ii Every $BCx=B(Cx)$ is in $\range B$. So $\range(BC)\subseteq\range B$ and $\rank(BC)\le\rank B$.
  \ii $\range(BC)=\{By:y\in\range C\}$ is the range of the restriction $B|_{\range C}$. By rank--nullity for this restriction, its dimension is $\dim\range C-\dim\nul(B|_{\range C})\le\rank C$.
  \ii (b) $B=C=\begin{psmallmatrix}0&1\\0&0\end{psmallmatrix}$: both have rank $1$, and $BC=0$ has rank $0$.
\end{steps}
\end{prob}

\begin{prob}\textbf{HW4 P5.}\par\idea{Put each basis in its own slot.}
\begin{steps}
  \ii Let $v_1,\dots,v_n$ and $w_1,\dots,w_m$ be bases of $V$ and $W$. Claim: $(v_1,0),\dots,(v_n,0),(0,w_1),\dots,(0,w_m)$ is a basis of $V\times W$.
  \ii Spans: $(v,w)=(\sum a_iv_i,\sum b_jw_j)=\sum a_i(v_i,0)+\sum b_j(0,w_j)$.
  \ii Independent: $\sum a_i(v_i,0)+\sum b_j(0,w_j)=(\sum a_iv_i,\sum b_jw_j)=(0,0)$ gives $\sum a_iv_i=0$ and $\sum b_jw_j=0$, so all $a_i=0$ and all $b_j=0$.
  \ii The list has $n+m$ vectors.
\end{steps}
\end{prob}

\begin{prob}\textbf{HW4 P6.}\par\idea{A subspace contains $0$; that locates $v$.}
\begin{steps}
  \ii (a) If $v+U$ is a subspace, $0\in v+U$, so $0=v+u$ and $v=-u\in U$.
  \ii If $v\in U$: $v+u\in U$ for all $u\in U$, so $v+U\subseteq U$; and each $u=v+(u-v)\in v+U$, so $U\subseteq v+U$. So $v+U=U$, a subspace.
  \ii (b) If $v+U=v'+U$, then $v=v+0\in v'+U$, so $v=v'+u$ and $v-v'\in U$.
  \ii If $v-v'\in U$: $v+u=v'+((v-v')+u)\in v'+U$, so $v+U\subseteq v'+U$; since $v'-v\in U$ as well, the reverse inclusion holds.
\end{steps}
\end{prob}

\begin{prob}\textbf{HW4 P7.}\par\idea{An invertible operator sends a basis to a basis. Build a basis of $V$ starting in $U$, and a second basis starting with the images under $S$.}
\begin{steps}
  \ii Let $u_1,\dots,u_m$ be a basis of $U$; extend to a basis $u_1,\dots,u_m,x_1,\dots,x_k$ of $V$.
  \ii $Su_1,\dots,Su_m$ is independent since $S$ is injective. Extend it to a basis $Su_1,\dots,Su_m,y_1,\dots,y_k$ of $V$ (the same $k$, since both bases have $\dim V$ vectors).
  \ii Define $T$ by $Tu_i=Su_i$, $Tx_j=y_j$. Then $T$ and $S$ agree on a basis of $U$, so $T|_U=S$.
  \ii $T$ sends a basis to a basis, so it is surjective, hence invertible (Prop.~31).
\end{steps}
\end{prob}

\begin{prob}\textbf{R1.}\par\idea{$S$ must undo $T$ on $\range T$; what it does elsewhere is free. Use a basis of $W$ that starts with $Tv_1,\dots,Tv_n$.}
\begin{steps}
  \ii (a) If $\sum\lambda_jTv_j=0$ then $T(\sum\lambda_jv_j)=0$, so $\sum\lambda_jv_j=0$, so all $\lambda_j=0$.
  \ii (b) Extend $Tv_1,\dots,Tv_n$ to a basis $Tv_1,\dots,Tv_n,w_1,\dots,w_k$ of $W$. Define $S$ by $S(Tv_j)=v_j$ and $S(w_i)=0$.
  \ii Then $ST(v_j)=v_j$ for each $j$: $ST$ and $I_V$ agree on a basis, so $ST=I_V$.
\end{steps}
\end{prob}

\begin{prob}\textbf{R2.}\par\idea{This is the map $\Gamma$.}
$\Gamma(u,w)=u+w$ is linear and maps onto $U+W$. If $u+w=0$ then $u=-w\in U\cap W=\{0\}$, so $u=w=0$: injective. So $\Gamma\colon U\times W\to U\oplus W$ is an isomorphism.
\end{prob}

\begin{prob}\textbf{R3.}\par\idea{Read off null space and range.}
(i) If $R(a)=0$ then every $a_j=0$: injective. The sequence $(1,0,0,\dots)$ is not in the range (every output starts with $0$): not surjective. (ii) For an operator on a finite-dimensional space, injective implies surjective (Prop.~31).
\end{prob}

\begin{prob}\textbf{R4.}\par\idea{Split $v$ as $(v-Pv)+Pv$.}
\begin{steps}
  \ii $Pv\in\range P$, and $P(v-Pv)=Pv-P^2v=0$, so $v-Pv\in\nul P$. Hence $V=\nul P+\range P$.
  \ii If $v\in\nul P\cap\range P$, write $v=Pw$. Then $v=Pw=P^2w=Pv=0$. So the sum is direct.
\end{steps}
\end{prob}

\begin{prob}\textbf{R5.}\par\idea{The matrix says $D\beta_1=\gamma_1$, $D\beta_2=\gamma_2$, $D\beta_3=\gamma_3$, $D\beta_4=0$. Put the null space last and let $\gamma$ be the images.}
$\beta=(x,\ x^2,\ x^3,\ 1)$ and $\gamma=(1,\ 2x,\ 3x^2)$. Both are bases (distinct degrees), and $Dx=1$, $Dx^2=2x$, $Dx^3=3x^2$, $D1=0$.
\end{prob}

\end{document}
